\BeleskeID{09_2-s030}
% element: 09_2-s030:e05
% element: 09_2-s030:e06
Iz uslova se vidi da postoji oblast u kojoj su oba tranzistora zasićena. Polazna ravnoteža u oznakama priključaka je
\[\frac{k_p}2\frac1{1+\frac{V_{GS,P}-V_{Tp}}{L_pE_{Cp}}}(V_{GS,P}-V_{Tp})^2=\frac{k_n}2\frac1{1+\frac{V_{GS,N}-V_{Tn}}{L_nE_{Cn}}}(V_{GS,N}-V_{Tn})^2\]
I sa modelom kratkog kanala, bez modulacije dužine kanala, ravnoteža ne sadrži izlazni napon. Zato se ponovo dobija beskonačno pojačanje. Tačka prebacivanja pripada ovoj oblasti. Sređivanje izraza obrađeno je u posebnom delu materijala „Podešavanje praga odlučivanja CMOS invertora“; ovo izvođenje često se pojavljuje na ispitu.

U realnosti modulacija dužine kanala vraća zavisnost od izlaznog napona, pa karakteristika nije vertikalna. Pojačanje i dalje može biti veoma veliko.
